\section{Equivalence at Matched Dimension}
\label{sec:equivalence}

The design uses $S_4$/$A_5$ as its contrast. The groups share $\dmin = 3$ but
lie on opposite sides of the solvability divide, the axis on which expressivity
theory separates state-tracking architectures
\citep{merrill2024illusion, grazzi2025negative,
siems2025deltaproduct}. If learned representational dimension were set
by complexity class, the pair should separate; if by representation
dimension, it should coincide. \citet{siems2025deltaproduct} provide the
precursor observation (\S5.2): $S_4$ and $A_5$ both extrapolate with two
Householder reflections and keys of size 3. They attribute this to the groups'
embedding in $\mathrm{SO}(3)$, and their PCA of the trained keys is
three-dimensional. The present test turns that single-architecture reading
into a pre-registered equivalence claim on state rank across the five-group
family. Section~\ref{sec:razor} makes the dimension causal. Because the
interesting outcome is a null, the pre-registration specifies an equivalence
test rather than a difference test. The comparison uses Welch two-one-sided
tests on restricted effective rank at margin $\pm 0.5$ rank-units (half the
spacing of the $\dmin$ ladder), with $n = 5$ seeds per group. The design record
fixed a pre-run power simulation before the sweep ran. It confirms that the
test reliably \emph{rejects} equivalence at a true gap of 1.0 rank-units, so a
real class effect of one ladder step could not have hidden inside the margin.

The observed difference is $+0.019$ rank-units (se 0.037, df 7.8), and
both one-sided tests pass at roughly seven times the critical value
($t = 13.06$ and $14.12$ against $t_{\mathrm{crit}} = 1.865$): the two
one-sided tests declare equivalence.
% <!-- evidence: N5 -->
The two groups' seed clusters are visually coincident in
Figure~\ref{fig:tracking} (lower panel). What the algebra demands governs
convergence here, not how hard the class is to compute.
